\chapter{Flow past an immersed cylinder: obstacles and forces}
\label{ch:D3_cylinder}
\noindent\texttt{tutorials/D\_flow/D3\_cylinder.py}\\[0.75em]
\begin{outcome}
You can compose everything: a carved exterior domain
(A6), the vector shifted-boundary condition for no-slip on the obstacle
(A3's machinery, per velocity component), a pseudo-time march to steady
flow (D2), and a FORCE observable — drag and lift from the shifted
traction on the surrogate boundary. You report a drag coefficient with
its configuration attached (confinement matters!).
\end{outcome}

Re = U D / nu = 20: steady, symmetric twin-vortex wake.
Domain [0,1]\textasciicircum{}2, cylinder r = 0.07 at (0.3, 0.5) — 14\% blockage, so expect
Cd ABOVE the unbounded-flow value (\textasciitilde{}2.0); our locked CI reference at
level 5 is Cd = 2.847. Forces come from
    F = oint ( p n\_hat - nu grad(u).n\_hat ) dS
on the surrogate staircase, area-corrected to the true circle, with
n\_hat the OBSTACLE-outward normal (orientation contract in
diffsim/sbm/vector.py — it was measured, not guessed).

\begin{expected}
\begin{verbatim}
Cd = 2.85 +/- 0.01,  |Cl| < 1e-3  (symmetric wake)
\end{verbatim}
\end{expected}

\section*{The script}
\begin{lstlisting}
import numpy as np
import scipy.sparse as sp
from scipy.sparse.linalg import splu

from diffsim.octree.build import build_uniform
from diffsim.mesh.nodes import build_mesh
from diffsim.mesh.constraints import build_constraints
from diffsim.mesh.basis import basis_tables
from diffsim.mesh.faces import face_tables
from diffsim.assembly.operators import DeviceMesh
from diffsim.geometry.csg import Sphere
from diffsim.sbm.surrogate import classify_lambda, extract_surrogate, GeometryData
from diffsim.sbm.vector import sbm_vector_dirichlet, surrogate_traction
from diffsim.api.ns_bricks import assemble_linear_ns
from diffsim.physics.poisson import gauss_points

DEVICE = "cuda:0"
R, CTR, U_IN = 0.07, (0.3, 0.5), 1.0
NU = 2 * U_IN * R / 20.0                        # Re_D = 20


if __name__ == "__main__":
    level, ndof, dim, dt = 5, 3, 2, 0.05
    # ---- carve the exterior domain and prepare the SBM data (A6 + A3) ----
    oracle = Sphere(CTR, R)
    tree = build_uniform(level, dim=2)
    ret, _ = classify_lambda(tree, oracle, 0.5, domain="outside")
    sf = extract_surrogate(ret)
    mesh = build_mesh(ret, p=1)
    cons = build_constraints(mesh)
    dm = DeviceMesh.from_mesh(mesh, cons, basis_tables(1, dim=2), DEVICE)
    geo = GeometryData.evaluate(oracle, ret, sf, face_tables(1, 2),
                                domain="outside")
    T = cons.T.tocsr()
    T_vec = sp.kron(T, sp.identity(ndof, format="csr"), format="csr")
    nfree = T.shape[1]
    coords = mesh.node_coords[cons.free_nodes]
    xq = gauss_points(mesh, dm.tables_by_p)
    on = lambda v, c: np.abs(coords[:, c] - v) < 1e-12
    strong = np.where(on(0.0, 0) | on(0.0, 1) | on(1.0, 1))[0]
    g_strong = np.zeros((len(strong), 2))
    g_strong[np.abs(coords[strong, 0]) < 1e-12, 0] = U_IN   # uniform inflow

    def gp_field(u_node):
        full = np.asarray(T @ u_node)
        aq, dq = {}, {}
        for pv in dm.bins:
            tb = dm.tables_by_p[pv]
            vals = full[mesh.conn_of[pv]]
            aq[pv] = np.einsum("qa,ead->eqd", tb.N, vals).reshape(-1, dim)
            h = mesh.tree.h()[mesh.bins[pv]]
            dq[pv] = (np.einsum("qad,ead->eq", tb.dN, vals)
                      * (2.0 / h)[:, None]).reshape(-1)
        return aq, dq

    # ---- pseudo-time march (D2's pattern, plus the SBM face block) ----
    x = np.zeros(nfree * ndof)
    sigma = 1.0 / dt
    qref = 0.5 * U_IN ** 2 * 2 * R
    for step in range(1, 161):
        u_node = x.reshape(nfree, ndof)[:, :dim]
        aq, dq = gp_field(u_node)
        fq = {pv: aq[pv] / dt for pv in xq}      # BDF1 history term
        A, b = assemble_linear_ns(dm, aq, dq, fq, NU, sigma=sigma)
        Af, bf = sbm_vector_dirichlet(          # no-slip on the cylinder
            dm, sf, geo, lambda y: np.zeros((len(y), 2)), NU, ndof)
        A = (A + T_vec.T @ Af @ T_vec).tolil()
        b = b + np.asarray(T_vec.T @ bf)
        for k, i in enumerate(strong):
            for c in range(dim):
                r = i * ndof + c
                A.rows[r] = [int(r)]
                A.data[r] = [1.0]
                b[r] = g_strong[k, c]
        pin = int(np.argmax(coords[:, 0] + coords[:, 1])) * ndof + dim
        A.rows[pin] = [pin]
        A.data[pin] = [1.0]
        b[pin] = 0.0
        x_new = splu(A.tocsr().tocsc()).solve(b)
        rate = np.abs(x_new - x).max() / dt
        x = x_new
        if step > 10 and rate < 5e-3:
            break
    F = surrogate_traction(dm, sf, geo, np.asarray(T_vec @ x), NU, ndof)
    print(f"steady after {step} steps:  Cd = {F[0] / qref:.3f}   "
          f"Cl = {F[1] / qref:+.5f}")
\end{lstlisting}

\begin{explore}
\begin{verbatim}
(a) Break the symmetry: move the cylinder to (0.3, 0.45). How large is
      Cl now, and which way does it point? Sanity-check the sign with the
      Bernoulli argument (faster flow on the wide side).
  (b) Raise Re to 60 (drop nu). The steady march stops converging — the
      wake wants to shed. Switch to a TRUE transient run (keep proper BDF2
      history, march ~10 convective times) and plot Cl(t): the von Karman
      street. Estimate the Strouhal number St = f D / U.
  (c) Blockage study: shrink r (and raise the level to keep D/h fixed).
      Cd should fall toward the unbounded ~2.0-2.1 literature band.
  (d) PERFORMANCE CORNER: the SBM face block (sbm_vector_dirichlet) is
      re-assembled every step although the geometry never moves. Cache it
      (it is additive!) and measure the saving. Then find the next
      biggest per-step cost and say what you would cache or fuse next.
\end{verbatim}
\end{explore}
